\begin{abstract}
Although large-scale unsupervised language models (LMs) acquire extensive world knowledge and exhibit some reasoning capabilities, exerting fine-grained control over their outputs remains challenging due to the absence of supervision in their training regimen. Standard approaches to address this challenge involve gathering human judgments that rank the quality of generated outputs and subsequently adjusting the LM via fine-tuning to match these human preferences—typically through reinforcement learning from human feedback (RLHF). Yet, RLHF introduces substantial complexity and instability by requiring a two-stage process: first, learning a reward model that captures human preferences, and then updating the LM through reinforcement learning to optimize this learned reward, all while maintaining similarity to the original LM. This work proposes a novel reward model parameterization for RLHF, enabling direct derivation of the optimal policy in closed form, which, in turn, transforms the RLHF objective into a straightforward classification loss. The proposed algorithm, termed \textit{Direct Preference Optimization} (DPO), is highly stable, efficient in terms of computational resources, and sidesteps the need for model sampling or extensive hyperparameter searches during fine-tuning. Our empirical evaluations reveal that DPO is capable of aligning LMs with human preferences as effectively as—or better than—existing techniques. Specifically, DPO outperforms RLHF methods based on PPO for sentiment control and achieves comparable or improved outcomes in summarization and single-turn dialogue, all while being considerably easier to implement and train.
\end{abstract}

\section{Introduction}
Large-scale unsupervised language models (LMs), when trained on extensive datasets, demonstrate a remarkable array of capabilities~\citep{chowdhery2022palm, brown2020language, touvron2023llama,bubeck2023sparks}. Despite their impressive proficiency, these models are exposed to human-generated data that reflects a broad spectrum of objectives, expertise, and priorities. Some of the strategies, intentions, or competencies present in this data are not ideal for the models to emulate. For instance, it is beneficial for an AI code assistant to \textit{recognize} prevalent programming mistakes to facilitate their correction, but it is preferable that, during code generation, the model emphasizes the (often infrequent) demonstrations of expert-level coding found within its training corpus. Similarly, while we expect the model to be \textit{informed} about widespread misconceptions held by a significant portion of people (e.g., 50\%), we would not want the model to propagate such inaccuracies in 50\% of responses to related questions. Thus, discerning and favoring the model’s \emph{targeted behaviors and outputs} from its extensive \textit{underlying capabilities and knowledge} is essential for constructing AI systems that are controllable, reliable, and aligned with user intentions \citep{ouyang2022training}. Conventional approaches generally guide LMs to reflect human value judgments through reinforcement learning (RL) techniques. Here, we will demonstrate that the prevailing RL-based preference optimization objective can be fulfilled precisely by a straightforward binary cross-entropy loss, considerably streamlining the overall process of preference alignment.

Broadly speaking, contemporary strategies incorporate human preferences, encoded in carefully designed datasets, to teach models the responses and behaviors regarded as helpful or safe. This alignment with human values is typically performed subsequent to the initial large-scale, unsupervised training on textual data. The simplest means of achieving this is supervised fine-tuning, where models are trained on examples of ideal behavior as provided by human annotators. Nonetheless, the most effective paradigm thus far is reinforcement learning from human (or AI-generated) feedback, known as RLHF/RLAIF \citep{christiano2017deep,bai2022constitutional}. RLHF involves training a reward model on preference data, then optimizing the LM’s policy through RL to maximize expected reward, while constraining divergence from the original pre-trained model. Although RLHF has produced models with remarkable performance in dialogue and programming tasks, its training pipeline is substantially more involved than supervised approaches. RLHF necessitates the training of multiple models and integrating online sampling of responses within the optimization loop, which results in high computational overhead.

In this study, we demonstrate a method for training language models to align with human preferences directly, circumventing the need for explicit reward modeling or reinforcement learning techniques. We introduce \textit{Direct Preference Optimization} (DPO), an approach that inherently targets the same optimization objective as traditional RLHF methods (maximizing reward under a KL-divergence penalty), while being notably easier to implement and train. Conceptually, the DPO method adjusts the model by enhancing the log odds of preferred outputs over non-preferred ones. It achieves this by applying an adaptive, instance-specific importance weighting, which helps prevent model collapse—a problem observed when employing a naive log-probability ratio as the objective. Like previous frameworks, DPO builds on a formal preference model (e.g., the Bradley-Terry model; \cite{bradley1952rankanalysis}) for quantifying how a reward function reflects observed human preferences. Distinctively, instead of employing the preference model solely to define and minimize a loss for training a separate reward model—followed by policy optimization using that model—DPO leverages a variable transformation, allowing the preference loss to be represented directly as a function of the policy itself. This enables the use of human preference datasets over model outputs to train the policy directly, using a simple binary cross-entropy loss, and to obtain the optimal policy with respect to a latent reward function derived from the preference annotations.

The primary innovation presented is Direct Preference Optimization (DPO): a straightforward, reinforcement learning-free strategy for optimizing language models with human preference data. Empirical evaluations indicate that DPO matches or surpasses the effectiveness of established approaches, such as RLHF with PPO, for preference-based tasks including sentiment adjustment, summarization, and conversational modeling, utilizing models containing up to 6B parameters.

\section{Related Work}

Self-supervised language models of growing size have demonstrated the capability to perform certain tasks without explicit training examples (zero-shot) \citep{radford2019language} or through the use of a handful of demonstrations within prompts (few-shot learning) \citep{gpt3,megatron,chowdhery2022palm}. Nonetheless, their effectiveness on various downstream tasks and alignment with user objectives can be markedly enhanced by applying fine-tuning on collections of instructional prompts paired with human-generated completions \citep{mishra-etal-2022-cross,sanh2022multitask,chung2022scaling,thoppilan2022lamda}. This process, referred to as ‘instruction-tuning’, equips large language models (LLMs) with the ability to extrapolate to previously unseen instructions and generally improves their practical utility \citep{chung2022scaling}. Even though instruction tuning has shown substantial success, collecting human preferences in the form of comparative judgments of model outputs is often more scalable than gathering expert demonstrations. Consequently, more recent research has focused on refining LLMs using datasets reflecting human choices, resulting in improved model competencies in translation \citep{kreutzer-etal-2018-reliability}, summarization \citep{stiennon2022learning,ziegler2020finetuning}, story generation \citep{ziegler2020finetuning}, and instruction-following tasks \citep{ouyang2022training,ramamurthy2023is}. These approaches typically begin by training a neural reward model to be consistent with observed preference data, often employing preference-based frameworks such as the Bradley-Terry model \citep{bradley1952rankanalysis}. Subsequently, the language model is adapted via reinforcement learning techniques to maximize these learned reward functions, utilizing algorithms like REINFORCE \citep{williams1992reinforce}, proximal policy optimization (PPO; \cite{schulman2017proximal}), or related variants \citep{ramamurthy2023is}. Closely aligned research directions make use of LLMs that have already been tuned with human feedback for instruction adherence to autonomously create further synthetic datasets of preferences, specifically targeting qualities like safety or harmlessness \citep{bai2022constitutional}. In these cases, weak supervision is provided through a rubric encoded as textual guidelines for LLM-generated annotations. Collectively, these techniques merge progress from two major research fronts: one focused on reinforcement learning-based training of language models for diverse aims~\citep{Ranzato2015SequenceLT,paulus2018a,wu2018learning}, and the other centered on general-purpose methods for learning from human preference data \citep{christiano2017deep,kupcsik2018learning}. Despite the benefits offered by relative preference optimization, effectively applying reinforcement learning for LLM fine-tuning remains a significant practical hurdle; the present study introduces a theoretically grounded strategy for optimizing over relative preferences that circumvents the use of reinforcement learning.

Beyond applications involving language, the study of learning policies from preferences has been explored in both bandit and reinforcement learning frameworks, resulting in various proposed methodologies. When actions are ranked or preferred rather than explicitly rewarded within a contextual bandit setup, the task is termed the contextual dueling bandit (CDB; \cite{yue2012karmed,dudik2015contextual}). In scenarios lacking explicit reward signals, CDBs employ the concept of a \textit{von Neumann winner} as a theoretical substitute for an optimal policy—defined as a policy achieving an expected win rate of at least 50\% when compared to \textit{any} alternative policy \citep{dudik2015contextual}. Notably, the CDB framework provides preference annotations in an online manner, whereas human preference learning usually operates on a pre-collected, static dataset of preference-labeled action pairs \citep{yan2022human}. Likewise, \textit{preference-based RL} (PbRL) leverages binary preferences produced by an \textit{unknown} scoring function instead of direct reward feedback \citep{BusaFekete2014,ruiz2023dueling}. Several algorithms address PbRL, some of which allow off-policy preference reuse, but the standard pipeline generally consists of an explicit estimation of the underlying scoring (reward) function before any policy optimization is performed \citep{jain2013learning,BusaFekete2014,christiano2017deep,sadigh2017active,kupcsik2018learning}. Contrasting these approaches, our work introduces a unified policy learning strategy that aims to optimize the policy in accordance with preferences in a single stage.

\section{Preliminaries}\label{section:prelims}

We summarize the RLHF methodology detailed by \citeauthor{ziegler2020finetuning}, which is further elaborated in subsequent works \citep{stiennon2022learning, bai2022training, ouyang2022training}. The pipeline typically unfolds in three stages: (1) supervised fine-tuning (SFT); (2) preference-based data collection with reward modeling; and (3) optimization through reinforcement learning.

\textbf{SFT}: RLHF pipelines customarily commence by taking a large language model and fine-tuning it on task-specific, high-quality data using supervised objectives, producing the model $\pi^\text{SFT}$ for the intended downstream applications (such as dialogue or summarization).

\textbf{Reward Modelling Phase}: Subsequently, this SFT model is prompted with inputs $x$, generating response pairs $(y_1, y_2)\sim \pi^\text{SFT}(y \mid x)$. Human annotators are then tasked with indicating a preferred response between the two, denoted by $y_w\succ y_l \mid x$, where $y_w$ and $y_l$ represent the favored and less favored responses, respectively, given the prompt $x$. These preferences are presumed to be governed by an unobserved reward function $r^*(y, x)$. Numerous models have been proposed to represent such preferences, with the Bradley-Terry (BT) \cite{bradley1952rankanalysis} model being a frequent selection; more general Plackett-Luce models \citep{plackett1975analysis, luce2012individual} are also applicable when dealing with multiple ranked outputs. According to the BT framework, the human preference probability distribution $p^*$ can be formulated as:

Given access to a fixed dataset of preference comparisons, $\mathcal{D}=\bigl\{x^{(i)}, y_w^{(i)}, y_l^{(i)}\bigr\}_{i=1}^N$, generated according to $p^*$, one can specify a reward function $r_{\phi}(x, y)$ with learnable parameters $\phi$ and fit it by maximizing the likelihood of observed pairwise preferences. Interpreting this setup as a binary classification problem, the negative log-likelihood objective can be formulated as follows:

\begin{equation}\label{eq:reward_model}
    \mathcal{L}_R(r_{\phi}, \mathcal{D}) = -\mathbb{E}_{(x, y_w, y_l)\sim \mathcal{D}}\bigl[\log \sigma(r_{\phi}(x, y_w)- r_{\phi}(x, y_l))\bigr]
\end{equation}

where $\sigma$ denotes the logistic sigmoid function. Within the large language model context, $r_{\phi}(x, y)$ typically builds upon the supervised fine-tuning (SFT) model $\pi^\text{SFT}(y \mid x)$, by adding a linear output head on top of the final transformer representation to predict a single reward scalar per sequence~\cite{ziegler2020finetuning}. To reduce reward estimate variance, previous studies enforce normalization so that $\mathbb{E}_{x,y\sim \mathcal{D}}\left[r_\phi(x, y)\right] = 0$ across $x$.

\textbf{RL Fine-Tuning Phase:} In this stage, the reward model is used to score language model outputs, providing the supervision signal. Consistent with prior research~\citep{jaques2017sequence, jaques2020human}, policy optimization aims to

\begin{equation}\label{eq:RL}
\max_{\pi_{\theta}}  \mathbb{E}_{x\sim \mathcal{D}, y\sim \pi_{\theta}(y \mid x)}\bigl[r_{\phi}(x, y)\bigr] - \beta\mathbb{D}_{\textrm{KL}}\bigl[\pi_{\theta}(y\mid x)\mid \mid \pi_\text{ref}(y\mid x)\bigr],
\end{equation}

where $\beta$ regulates how much the current policy $\pi_\theta$ is permitted to diverge from a reference, generally the initial SFT model $\pi_\text{ref} = \pi^\text{SFT}$. In practice, training starts from $\pi_\theta = \pi^\text{SFT}$ as well. This KL regularization acts as an anchor to keep the policy close to data regimes where the reward model is reliable, encourages response diversity, and counters degeneracy such as mode collapse. Since language generation involves discrete outputs, the objective is not directly differentiable and must be tackled using reinforcement learning. Standard approaches~\citep{ziegler2020finetuning, stiennon2022learning, bai2022training, ouyang2022training} define a combined reward ${r(x, y) = r_{\phi}(x, y) -\beta (\log \pi_{\theta}(y\mid x) - \log \pi_\text{ref}(y\mid x))}$, and optimize it with algorithms such as PPO~\cite{schulman2017proximal}.

\section{Direct Preference Optimization}\label{sec:DPO}

Recognizing the complexities and scalability issues of reinforcement learning for large models, we propose a streamlined method to directly optimize policies from preference data, thereby bypassing explicit reward maximization via RL. Unlike established RLHF pipelines that independently fit a reward model before conducting RL-based policy improvement, our framework selects a reward parametrization which admits a closed-form solution for the optimal policy, thus eliminating the need for an iterative RL procedure. Central to this method is the exploitation of an exact, analytical correspondence between reward models and their optimal policies; this facilitates a reparameterization that recasts loss minimization over rewards into policy space. In doing so, we avoid training an isolated reward predictor, but still operate within preference modeling structures such as the Bradley-Terry framework. As a result, our policy model serves a dual role—simultaneously representing the language model and embodying an implicit reward function.

\textbf{Derivation of the DPO Objective.} Our starting point is the standard reinforcement learning objective with a generic reward function $r$, as expressed in Eq.~\ref{eq:RL}, consistent with previous literature. Building upon earlier studies~\citep{peters2007reinforcement, peng2019advantage, korbak2022reinforcement, go2023aligning}, it can be readily established that the solution maximizing the KL-constrained objective from Eq.~\ref{eq:RL} is characterized by:

\begin{equation}\label{eq:op_policy}
    \pi_r(y\mid x) = \frac{1}{Z(x)}\pi_\text{ref}(y\mid x)\exp\left(\frac{1}{\beta}r(x, y)\right),
\end{equation}

with $Z(x) =\sum_{y}\pi_\text{ref}(y\mid x)\exp\left(\frac{1}{\beta}r(x, y)\right)$ serving as the normalization constant. The complete derivation appears in Appendix \ref{app:derivation1}. Despite using the maximum likelihood estimate $r_{\phi}$ as a surrogate for the true reward $r^*$, direct computation or approximation of $Z(x)$ remains challenging and computationally prohibitive~\citep{korbak2022reinforcement, go2023aligning}, making direct utilization of this formulation impractical. Nevertheless, Eq.~\ref{eq:op_policy} can be re-expressed so that the reward depends on the policy $\pi_r$, the reference distribution $\pi_\text{ref}$, and the (unknown) partition function $Z(\cdot)$. To achieve this, we first apply the logarithm to both sides of Eq.~\ref{eq:op_policy} and rearrange:

\begin{equation}\label{eq:main_eq}
    r(x,y) =\beta \log \frac{\pi_r(y\mid x)}{\pi_\text{ref}(y\mid x)} + \beta \log Z(x).
\end{equation}

This reparameterization is valid for both the ground-truth reward $r^*$ and the associated optimal policy $\pi^*$. Importantly, the Bradley-Terry model only involves the reward difference between two outputs: ${p^*(y_1 \succ y_2 \mid x) = \sigma(r^*(x, y_1) - r^*(x, y_2))}$. If we substitute $r^*(x,y)$ using Eq.~\ref{eq:main_eq} into the preference model Eq.~\ref{eq:bradley-terry}, the terms involving the partition function cancel out, resulting in a preference probability expression depending solely on the optimal policy $\pi^*$ and the reference policy $\pi_\text{ref}$. Consequently, the optimal RLHF policy $\pi^*$ within the Bradley-Terry model framework adheres to:

\begin{equation}\label{eq:objective}
    p^*(y_1\succ y_2 \mid x)=\frac{1}{1 + \exp\left(\beta \log \frac{\pi^*(y_2\mid x)}{\pi_\text{ref}(y_2\mid x)} - \beta \log \frac{\pi^*(y_1\mid x)}{\pi_\text{ref}(y_1\mid x)}\right)}
\end{equation}

A detailed derivation appears in Appendix~\ref{app:derivation2}. Although Eq.~\ref{eq:objective} specifically adopts the Bradley-Terry formulation, similar logic extends to generalizations such as the Plackett-Luce models~\citep{plackett1975analysis, luce2012individual}, discussed in Appendix~\ref{app:plackett_luce_models}.

Having linked the human preference probabilities directly to the optimal policy instead of the reward model, it becomes feasible to optimize a maximum likelihood objective over a parameterized policy $\pi_\theta$. Mirroring the procedure in reward modeling (see Eq.~\ref{eq:reward_model}), the resulting policy learning objective is:

\begin{equation}\label{eq:optimum_model}
    \mathcal{L}_\text{DPO}(\pi_{\theta}; \pi_\text{ref}) = -\mathbb{E}_{(x, y_w, y_l)\sim \mathcal{D}}\left[\log \sigma \left(\beta \log \frac{\pi_{\theta}(y_w\mid x)}{\pi_\text{ref}(y_w\mid x)} - \beta \log \frac{\pi_{\theta}(y_l\mid

In this approach, we infer an implicit reward through an alternative parameterization, ensuring that its optimal policy aligns directly with $\pi_\theta$. Additionally, as our methodology amounts to reparameterizing the Bradley-Terry model, it inherits theoretical guarantees such as consistency under appropriate conditions on the distribution of preference data \cite{bong2022generalized}. We examine these theoretical attributes of DPO in further depth in Section~\ref{sec:theory}, situating them within the broader literature.

\textbf{How does the DPO update function?} To gain a mechanistic perspective on DPO, it is instructive to study the gradient of the loss $\mathcal{L}_\text{DPO}$. The derivative of this objective with respect to the model parameters $\theta$ can be expressed as:

\begin{multline*}\label{eq:gradient}
    \nabla_\theta \mathcal{L}_\text{DPO}(\pi_\theta;\pi_\text{ref}) = \\ -\beta\mathbb{E}_{(x, y_w, y_l) \sim \mathcal{D}} \bigg[\underbrace{\sigma(\hat{r}_\theta(x, y_l) - \hat{r}_\theta (x, y_w))}_\text{assigns higher emphasis if reward order is mistaken}\bigg[\underbrace{\nabla_\theta\log \pi(y_w \mid x)}_\text{boosts probability of $y_w$} - \underbrace{\nabla_\theta\log\pi(y_l \mid x)}_\text{reduces probability of $y_l$}\bigg]\bigg],
\end{multline*}

where the implicit reward is specified by $\hat{r}_\theta(x, y) = \beta \log \frac{\pi_\theta(y \mid x)}{\pi_\text{ref}(y \mid x)}$, determined by both the language model $\pi_\theta$ and the reference model $\pi_\text{ref}$ (with further discussion in Section~\ref{sec:theory}). Conceptually, this gradient step increases the probability assigned to favorable completions $y_w$ and suppresses the probability of unfavorable ones $y_l$. Critically, the influence of each data point depends on the extent to which the implicit reward function $\hat{r}_\theta$ erroneously favors dispreferred answers, with the scale governed by $\beta$, effectively quantifying the severity of the ranking mistake and incorporating the KL divergence constraint. Our empirical analysis reveals this scaling is essential—simplified variants that omit this weighting factor lead the language model toward pathological behaviors (see Appendix Table~\ref{tab:unlikelihood_generations}).

\textbf{DPO overview.} The typical workflow for DPO comprises: 1) Generating completions $y_1, y_2 \sim \pi_\text{ref}(\cdot \mid x)$ for each input prompt $x$, then assigning human preference labels to form the offline preference dataset $\mathcal{D} = \{x^{(i)}, y_w^{(i)}, y_l^{(i)}\}_{i=1}^N$; 2) training the language model $\pi_\theta$ by minimizing the DPO loss $\mathcal{L}_\text{DPO}$, using the selected $\pi_\text{ref}$, dataset $\mathcal{D}$, and the target $\beta$. In practical scenarios, it is preferable to utilize existing publicly available preference datasets instead of creating new samples and collecting fresh human annotations. As these datasets are usually created using $\pi^\text{SFT}$, we set $\pi_\text{ref} = \pi^\text{SFT}$ when possible. Conversely, if $\pi^\text{SFT}$ is unavailable, we define $\pi_\text{ref}$ by maximizing the likelihood of the favored completions ${(x, y_w)}$, specifically, ${\pi_\text{ref} = \argmax_{\pi}\mathbb{E}_{x, y_w \sim \mathcal{D}}\left[\log \pi(y_w \mid x)\right]}$. This method reduces the discrepancy between the inaccessible true reference distribution and the surrogate $\pi_\text{ref}$ employed in DPO. More comprehensive details on the implementation specifics and hyperparameter selections are included in Appendix~\ref{app:implementation}.

\section{Theoretical Analysis of DPO} In this section, we further elucidate the DPO approach, present theoretical justification, and connect the strengths of DPO to common limitations found in actor-critic algorithms (e.g., PPO~\cite{schulman2017proximal}) that are utilized for RLHF.

\label{sec:theory}

\subsection{Your Language Model Is Secretly a Reward Model} DPO avoids both learning an explicit reward function and performing reinforcement learning to train the policy, instead leveraging a single maximum likelihood objective. Observe that the objective in Eq.~\ref{eq:main_eq} coincides with a Bradley-Terry model whose reward is parameterized as $r^*(x, y) = \beta \log\frac{\pi^*_\theta(y \mid x)}{\pi_\text{ref}(y \mid x)}$, and the parametric model $\pi_{\theta}$ is optimized analogously to reward model learning as formalized in Eq.~\ref{eq:reward_model}, after an appropriate variable substitution. Here, we will develop the theoretical basis for this reparameterization, demonstrate that it does not limit the class of reward models that can be learned, and show that it supports the exact recovery of the optimal policy. We begin by introducing an equivalence relation for reward functions.

\begin{definition}
We define two reward functions $r(x, y)$ and $r'(x, y)$ to be equivalent if and only if ${r(x, y)-r'(x, y) = f(x)}$ for some function $f$.
\end{definition}

This notion of equivalence constitutes an equivalence relation, thereby dividing the set of reward functions into equivalence classes. We present two lemmas as follows:

\begin{lemma}\label{lemma:same_prefrence}
Within the Plackett-Luce framework—including, as a special case, the Bradley-Terry model—any two reward functions belonging to the same equivalence class will yield identical distributions over preferences.
\end{lemma}

\begin{lemma}\label{lemma:same_policy}
    Any pair of reward functions within the same equivalence class produces the same optimal policy in the associated constrained RL problem.
\end{lemma}

Proofs for these statements are relatively simple and can be found in Appendix~\ref{app:lemma1}. The first lemma highlights a well-known identification problem inherent to the Plackett-Luce class of models~\cite{plackett1975analysis}: the inability to uniquely recover a reward function without further constraints. Consequently, identifiability conditions must be imposed to obtain meaningful maximum likelihood estimates as per Eq.~\ref{eq:reward_model}~\cite{bong2022generalized}. The second lemma asserts that every member of an equivalence class leads to the same optimal policy; therefore, in the context of our objectives, it suffices to recover any representative reward from the correct equivalence class. The subsequent theorem is established in Appendix~\ref{app:thm1}:

\begin{theorem}\label{thm:main}
    Subject to mild technical assumptions, any equivalence class of rewards compatible with the Plackett-Luce (and, in particular, the Bradley-Terry) model can be expressed as ${r(x, y) = \beta \log \frac{\pi(y\mid x)}{\pi_\text{ref}(y\mid x)}}$ for a certain model $\pi(y\mid x)$ and a fixed reference model $\pi_\text{ref}(y \mid x)$.
\end{theorem}

\begin{sproof}
    Suppose we are given any reward function $r(x, y)$, which corresponds to an optimal policy $\pi_r(y \mid x)$ as defined in Eq.~\ref{eq:op_policy}. We aim to demonstrate that some member of the equivalence class containing $r$ admits the above reparameterization. Let us introduce the following projection operator:
\begin{equation}
    f(r; \pi_\text{ref}, \beta)(x, y) = r(x, y) - \beta\log\sum_{y}\pi_\text{ref}(y\mid x)\exp\left(\frac{1}{\beta}r(x, y)\right)
\end{equation}
This operation amounts to normalizing the original reward by subtracting the log-partition function, which depends solely on $x$. Therefore, $f(r; \pi_\text{ref}, \beta)(x, y)$ belongs to the same equivalence class as $r(x, y)$. Substituting the right-hand side of Eq.~\ref{eq:main_eq} (which universally holds), we deduce $f(r; \pi_\text{ref}, \beta)(x, y) = \beta \log \frac{\pi_r(y\mid x)}{\pi_\text{ref}(y\mid x)}$. Thus, the projection operator yields an equivalence-class representative in the desired parametric form, establishing that our proposed reparameterization retains generality for modeling rewards.
\end{sproof}

Alternatively, Theorem~\ref{thm:main} can be interpreted as identifying the specific reward function, within each equivalence class, that the DPO reparameterization chooses—namely, the reward function that satisfies

\begin{equation}\label{eq:lag_p}
     \sum_{y}\underbrace{\pi_\text{ref}(y\mid x)\exp\left(\frac{1}{\beta}r(x, y)\right)}_{=\pi(y\mid x)\text{, via Thm.~\ref{thm:main} reparam.}} = 1,
\end{equation}

which ensures that $\pi(y\mid x)$ constitutes a well-defined probability distribution (i.e., non-negative and summing to one). Notably, as per Eq.~\ref{eq:op_policy}, this normalization in Eq.~\ref{eq:lag_p} functions as the partition function for the optimal policy derived from the reward function $r(x, y)$. The central observation behind the DPO algorithm is the ability to impose suitable constraints on the Plackett-Luce (and particularly Bradley-Terry) preference model families—which are otherwise under-determined—while maintaining the expressivity of the corresponding reward models, and at the same time rendering the optimal policy in Eq.~\ref{eq:op_policy} analytically computable for every prompt $x$.

\subsection{Instability of Actor-Critic Algorithms}
The presented framework also allows us to analyze and better understand the instabilities observed in common actor-critic approaches for RLHF, including PPO. Adhering to the RLHF pipeline, we restrict attention to the RL fine-tuning procedure as detailed in Section \ref{section:prelims}. This discussion draws parallels to the control as inference framework \cite{levine2018reinforcement} when applied to the constrained RL scenario described by \ref{eq:RL}. We suppose a parameterized distribution $\pi_{\theta}(y\mid x)$ and aim to minimize $\mathbb{D}_{\text{KL}}[\pi_{\theta}(y|x) \mid \mid \pi^*(y\mid x)]$, with $\pi^*$ denoting the optimal policy, as defined in Eq.~\ref{eq:optimum_model}, corresponding to the reward $r_{\phi}(y, x)$. By expanding the relevant terms, the resultant objective takes the form

\begin{equation}\label{eq:AC}
    \max_{\pi_{\theta}}\mathbb{E}_{\pi_{\theta}(y\mid x)}\left[\underbrace{r_{\phi}(x, y) -\beta\log\sum_{y}\pi_\text{ref}(y\mid x)\exp\left(\frac{1}{\beta}r_{\phi}(x, y)\right)}_{f(r_{\phi}, \pi_\text{ref}, \beta)} - \underbrace{\beta\log\frac{\pi_{\theta}(y\mid x)}{\pi_\text{ref}(y\mid x)}}_{\text{KL}}\right]
\end{equation}

The objective considered here aligns with the formulation employed in earlier studies 
\citep{ziegler2020finetuning, stiennon2022learning, bai2022training, ouyang2022training}, which utilize the DPO-equivalent reward within the $r_{\phi}$ reward class. Within this framework, the normalization factor present in $f(r_{\phi}, \pi_\text{ref}, \beta)$ can be interpreted as the soft value function of the reference policy $\pi_\text{ref}$. Although omitting this term leaves the optimal policy unchanged, its absence can significantly increase the variance of policy gradients and thereby destabilize the training process. Addressing the normalization term through a value function estimator is possible but may lead to optimization challenges. Traditionally, normalization has been handled in previous work by employing a human completion baseline—a Monte Carlo estimate with a single sample—for the normalization constant. Notably, the DPO reparameterization leads to a reward formulation that eliminates the need for such baselines altogether.

\section{Experiments}
Here, we conduct empirical studies to assess how effectively DPO can train policies based solely on preference data. Our initial experiments utilize a controlled text generation scenario to analyze DPO’s balance between maximizing expected reward and controlling KL-divergence with respect to the reference policy, relative to established preference learning approaches like PPO. Subsequent experiments extend this analysis to more complex tasks and larger models, focusing on RLHF applications such as summarization and conversational dialogue. Across these tasks, we observe that DPO typically matches or surpasses strong baselines, including PPO-based RLHF and strategies that select the top trajectory from $N$ samples under a learned reward, with minimal hyperparameter tuning. A detailed description of our experimental methodology precedes these results; supplementary information can be found in Appendix~\ref{app:exp_details}.

\textbf{Tasks.} In our study, we investigate three distinct tasks focused on open-ended text generation. Throughout all experiments, each method is tasked with learning a policy from a dataset of human preferences denoted as $\mathcal{D}=\bigl\{x^{(i)}, y_w^{(i)}, y_l^{(i)}\bigr\}_{i=1}^N$. For the \textbf{controlled sentiment generation} task, $x$ corresponds to an initial segment of a movie review from the IMDb dataset \cite{maas-EtAl:2011:ACL-HLT2011}, and the goal is for the policy to generate a continuation $y$ that expresses positive sentiment. To enable a systematic assessment, we \textit{construct} pairs of preferences by using an already trained sentiment classifier to score generated continuations, accepting only those pairs for which $p(\text{positive}\mid x,y_w)>p(\text{positive}\mid x,y_l)$. For SFT, GPT-2-large is fine-tuned to convergence using the training portion of the IMDB dataset (additional experimental protocol in App~\ref{app:sentiment_details}). The \textbf{summarization} task involves $x$ as a Reddit forum post, where the aim is to produce a summary $y$ capturing its key messages. Consistent with existing research, we make use of the Reddit TL;DR summarization dataset \citep{volske-etal-2017-tl} and human preference data assembled by \citeauthor{stiennon2022learning}. The SFT model in this task is trained on a dataset of human-composed summaries of forum posts\footnote{\url{https://huggingface.co/CarperAI/openai_summarize_tldr_sft}}, leveraging the TRLX \citep{leandro_von_werra_2023_7790115} implementation for RLHF. Notably, the human feedback data was annotated on outputs of a separate but similarly tuned SFT model, as documented by \citeauthor{stiennon2022learning}. For the third task, \textbf{single-turn dialogue}, $x$ is a user-issued prompt, which may range from factual questions to personal requests. Here, the policy must generate a relevant and constructive reply $y$; we adopt the Anthropic Helpful and Harmless dialogue dataset \citep{bai2022training}, consisting of 170,000 exchanges between a user and an AI assistant. Each interaction concludes with two AI-generated replies and an annotation indicating which one the human preferred. Unlike the other tasks, an off-the-shelf SFT model is not available in this context, so we create one by fine-tuning a pretrained language model exclusively on preferred responses.

\textbf{Evaluation.} We adopt two distinct strategies for evaluation in our experimental analysis. To systematically assess each algorithm’s capacity to optimize the objective of constrained reward maximization, we conduct controlled sentiment generation experiments where each method is evaluated according to the frontier that trades off realized reward and KL-divergence from the reference policy; this evaluation is feasible since we possess the exact reward function (a sentiment classifier). Conversely, outside of controlled settings—where the ground-truth reward function is unavailable—algorithms are compared using their \textit{win rate} over a baseline policy, with GPT-4 serving as a surrogate for human judges to score both summary quality (in summarization) and helpfulness (in single-turn dialogue). In summarization, the test set’s gold reference summaries serve as the baseline; for dialogue, the baseline is the annotated preferred response from the test set. Although prior work indicates that LMs can outperform traditional metrics as automatic evaluators \citep{Chen2023ExploringTU}, we substantiate our reliance on GPT-4 via a dedicated human evaluation, described in Sec.~\ref{sec:human-judgments}. Our findings indicate that human raters show strong concordance with GPT-4, with human-GPT-4 agreement matching or exceeding agreement among human annotators.

\textbf{Methods.} Alongside DPO, we benchmark a variety of established approaches for aligning language models with human preferences. For the summarization task, we employ zero-shot prompting of \textbf{GPT-J} \citep{gpt-j}, while for the dialogue setting we utilize 2-shot prompting with \textbf{Pythia-2.8B} \citep{biderman2023pythia}. We also include evaluations of the \textbf{SFT} model, as well as \textbf{Preferred-FT}—which is produced via supervised fine-tuning on the selected completions $y_w$, using either the SFT output (in sentiment and summarization) or that of a general language model (in dialogue). Another relevant technique, \textbf{Unlikelihood}~\citep{welleck2019neural}, optimizes for higher probability on $y_w$ while directly penalizing likelihood of $y_l$; here, an optional weighting coefficient $\alpha\in[0,1]$ can be assigned to the unlikelihood penalty term. Moreover, our comparison set includes \textbf{PPO} \citep{schulman2017proximal}, utilizing rewards inferred from preference data, and \textbf{PPO-GT}, which leverages the oracle access to the exact reward function available in the sentiment-controlled context. In the sentiment experiments, two versions of PPO-GT are examined: an unmodified implementation from \cite{leandro_von_werra_2023_7790115}, and an adapted version that normalizes rewards and fine-tunes hyperparameters for enhanced outcomes (these enhancements are similarly applied to PPO runs using learned reward models). As another point of reference, we include the \textbf{Best of $N$} strategy: this involves sampling $N$ candidate outputs from the SFT model (or Preferred-FT for dialogue), then selecting the response with the highest score under a reward model trained from preference data. This approach can yield high-performing outputs while disentangling reward modeling from PPO’s optimization, but its computational costs scale poorly with $N$ since it requires generating $N$ completions per query at test time.

\subsection{How well can DPO optimize the RLHF objective?}

RLHF algorithms typically employ a KL-constrained reward maximization objective, which seeks to maximize reward while simultaneously preventing the policy from straying too far from the reference policy. Consequently, a meaningful comparison of different methods necessitates examining both achieved reward and the associated KL divergence; obtaining a marginally higher reward is insufficient if it comes at the expense of substantially increased KL divergence. Figure~\ref{fig:frontier-tldr-main} depicts the tradeoff frontier between reward and KL for various methods in the sentiment task. For each method, we conduct several independent training runs, each configured with a distinct policy conservativeness parameter (using target KL values $\in\{3,6,9,12\}$ for PPO, $\beta \in \{0.05,0.1,1,5\}$, and $\alpha\in\{0.05,0.1,0.5,1\}$ for unlikelihood; preferred-FT is varied via random seed). In total, this sweep covers 22 experimental runs. At intervals of 100 training steps until convergence, policies are evaluated on a test set of prompts, calculating both the mean reward according to the actual reward function and the mean sequence-level KL\footnote{Specifically, the aggregate KL-divergence over all time steps in a sequence.} relative to the reference policy $\text{KL}\left(\pi\mid \mid \pi_\text{ref}\right)$. The results demonstrate that DPO consistently yields the most favorable reward-KL frontier, attaining the highest reward while maintaining a low KL. This finding is especially remarkable for several reasons. To begin with, even though DPO and PPO optimize the identical objective, DPO shows significantly greater efficiency—its reward/KL tradeoff is strictly better than that of PPO. Furthermore, DPO outperforms PPO not only under standard conditions, but also \emph{when PPO is supplied with access to ground truth rewards} (PPO-GT).

\subsection{Can DPO scale to real preference datasets?}
\label{sec:dpo-real-datasets}
We subsequently assess how well DPO can fine-tune language models on tasks involving summarization and single-turn dialogue, using real human preference data. In the case of summarization, widely-used automatic metrics like ROUGE are often weak proxies for actual human judgments~\citep{stiennon2022learning}. Previous research indicates that leveraging PPO for LM fine-tuning based on human preference signals leads to more effective summary outputs. To compare various approaches, we generate outputs for the test set of the TL;DR summarization dataset, and determine each method’s average win rate relative to reference summaries. Generation occurs across a temperature range from 0.0 to 1.0, and the resulting win rates are visualized in Figure~\ref{fig:frontier-tldr-main} (right). The DPO, PPO, and Preferred-FT methods all adapt the same GPT-J SFT model\footnote{\url{https://huggingface.co/CarperAI/openai_summarize_tldr_sft}}. Experimental results reveal that DPO attains a win rate close to 61\% at temperature 0.0, outperforming PPO’s peak of approximately 57\% at its optimal temperature. Additionally, DPO demonstrates a greater maximal win rate than the best-of-$N$ approach. It is important to highlight that DPO’s $\beta$ hyperparameter received minimal tuning in these experiments, suggesting that further performance gains might be achievable. Notably, DPO maintains its effectiveness over a broader temperature span than PPO, which, at elevated temperatures, performs comparably to the original GPT-J model. Preferred-FT, by contrast, yields little improvement relative to the SFT baseline. In direct human evaluation comparisons presented in Section~\ref{sec:human-judgments}, outputs generated by DPO at temperature 0.25 are preferred 58\% of the time when compared to PPO outputs at the same temperature.

For single-turn dialogue evaluation, we focus on the portion of the Anthropic HH test set \citep{bai2022training} that contains only one exchange between a human and an assistant. To assess the win rates of various approaches, we use GPT-4 to judge responses by comparing each method’s outputs against the human-preferred completions within the test set. Since a conventional SFT model for this benchmark does not exist, we begin with a pre-trained Pythia-2.8B model, employ Preferred-FT to build a reference model trained on the preferred completions to ensure output alignment with the model’s distribution, and subsequently fine-tune using DPO. In addition, we benchmark against two baselines: the strongest sample among 128 Preferred-FT completions (noting that performance for the Best of $N$ strategy saturates at $N=128$ for this dataset; refer to Appendix Figure~\ref{fig:best-of-n}) and a two-shot prompt setup with the Pythia-2.8B base model, observing that DPO matches or outperforms other techniques at their optimal temperature settings. Furthermore, we analyze a PPO-based RLHF model trained on the Anthropic HH data \footnote{\url{https://huggingface.co/reciprocate/ppo_hh_pythia-6B}} sourced from a reputable implementation \footnote{\url{https://github.com/CarperAI/trlx/tree/main/examples/hh}}; however, no prompt or sampling temperature configuration yields results surpassing those of the original Pythia-2.8B model. Drawing on findings from TL;DR and the fact that PPO and Best of 128 both target the same reward function, we regard the Best of 128 method as an approximate indicator of PPO-level success. Ultimately, DPO stands out as the only method that both surpasses the preferred completions from the Anthropic HH dataset and achieves parity or superiority relative to the resource-intensive Best of 128 approach, all while maintaining computational efficiency. As highlighted in Figure~\ref{fig:dialogue-main}, DPO also attains its optimal performance in a comparatively small number of steps.

\subsection{Generalization to a new input distribution}

To more thoroughly assess how PPO and DPO perform under distribution shifts, we examine the behavior of policies trained in our Reddit TL;DR summarization experiment by testing them on an alternative domain: the test split of the CNN/DailyMail news article dataset \citep{nallapati-etal-2016-abstractive}. We apply the optimal sampling temperatures identified from the TL;DR experiments (0 and 0.25) in this evaluation. Table~\ref{tab:ood} summarizes the outcomes. As in the previous setup, we calculate the GPT-4 win rate over ground-truth summaries, utilizing the GPT-4 (C) prompt with a minor modification—substituting “forum post” for “news article.” Notably, on this out-of-domain dataset, DPO still surpasses PPO by a clear margin. These findings offer preliminary support that DPO achieves at least comparable generalization to PPO, even without access to the extra unlabeled Reddit TL;DR prompts available to PPO.

\subsection{Validating GPT-4 judgments with human judgments}
\label{sec:human-judgments}
To assess the trustworthiness of GPT-4’s comparative evaluations, we conduct a human subject study using outputs from our TL;DR summarization experiments and two different GPT-4 prompt formulations. The \textbf{GPT-4 (S)} (“simple”) prompt asks directly which summary best captures the salient information in a post. The \textbf{GPT-4 (C)} (“concise”) prompt, in addition, requests a decision about which summary is more concise—this variant is included because, using the simple prompt, GPT-4 tends to select longer, more redundant summaries than humans do. The full text of each prompt is provided in Appendix~\ref{app:prompts}. In our study, we select three representative system configurations for comparison: the best DPO setting (temperature 0.25), the weakest PPO setting (temperature 1.0), and an intermediate SFT setting (temperature 0.25). Each is contrasted against the best-performing greedy PPO baseline. This arrangement ensures a broad spectrum of summary quality is included. Analysis reveals that, across both prompt types, the rate at which GPT-4’s judgments coincide with human assessments is similar to the agreement rate among human annotators themselves—implying GPT-4 serves as a suitable proxy for human evaluation (though for DPO and PPO-1, multiple independent human judgments are gathered due to rater availability constraints). Furthermore, we observe that the \textbf{GPT-4 (C)} prompt’s win rates more closely mirror human preferences, prompting us to adopt it for the principal experiments in Section~\ref{sec:dpo-real-datasets}. Further details on the human evaluation protocol, interface, and volunteer raters are provided in Appendix~\ref{app:human-study}.

\section{Discussion}
Preference-based learning offers a robust and scalable approach for developing capable and aligned language models. In this work, we proposed DPO, a straightforward training method that enables language models to be optimized based on preference data without the use of reinforcement learning. Rather than forcing preference learning to fit into a conventional RL framework and relying on standard RL algorithms, DPO establishes a correspondence between language model policies and reward functions, permitting direct optimization according to human preferences using a basic cross-entropy loss. This method bypasses both the need for reinforcement learning and any loss in generality. Despite requiring virtually no hyperparameter tuning, DPO matches or outperforms leading RLHF methods, such as those employing PPO, thereby lowering the barriers to leveraging human preferences for training language models.

\textbf{Limitations \& Future Work.} Our findings prompt several key questions for subsequent investigation. What is the out-of-distribution generalization capability of DPO-trained policies compared to methods that rely on an explicit reward signal? Preliminary evidence indicates that DPO can generalize in a manner comparable to PPO-based models, yet further systematic analysis is necessary. For instance, might employing self-labeling with the DPO policy enable effective use of prompts without labels? Another area of inquiry concerns how reward over-optimization arises in the DPO framework, and whether the modest decline in performance observed in Figure~\ref{fig:dialogue-main}-right may be attributable to this phenomenon. Moreover, although our current evaluation is limited to models with up to 6B parameters, future work could extend DPO to cutting-edge models with much larger scales. In terms of evaluation protocols, our observations show that GPT-4 win rates are influenced by the choice of prompt, suggesting a need for further research into eliciting reliable judgments from automated evaluators. Ultimately, DPO may have valuable uses well beyond language model preference alignment, including its potential application for training generative models across a range of data modalities.